\section{A quadratic jet bound on diagonal products}\label{sec:diagonal}

The scalar estimates control section orders on a canonical variety and
local positivity on its index cover. We next obtain a complementary
upper bound on products of that cover. The bound grows quadratically
with the number of factors and is independent of the covering degree.
We give the arbitrary-dimensional form of the diagonal argument in
\cite[Section~7]{U4}, including its positive-characteristic estimates.

For a divisor $L$ on a variety $V$, write
$L^{\boxplus t}=\sum_{i=1}^t\operatorname{pr}_i^*L$ on $V^t$.
The order invariant $\gamma$ and the very general Seshadri constant
$\varepsilon$ have the conventions used in Proposition~\ref{prop:scalar}.

\begin{proposition}[Diagonal jet bound]\label{prop:diagonal}
Let $V$ be a normal projective canonical complex variety of dimension
$n\ge2$, with $K_V\sim_{\Q}0$ and canonical index $r$. Let
$\pi:Y\to V$ be its canonical cyclic cover. Suppose that $L$ is an
ample rational Cartier divisor and $C_1>0$ is rational, with
\[
 1\le L^n\le2,\qquad \gamma(L;V)\le C_1.
\]
Choose a positive rational number $b_0$ such that
\begin{equation}\label{eq:diag-b0}
 3(2b_0)^n<\frac14,\qquad b_0(C_1+1)<\frac34.
\end{equation}
For every integer $t\ge2$, set
\[
 Z_t=Y^t/\mu_{r,\mathrm{diag}},\qquad
 \theta_t:Z_t\longrightarrow V^t,\qquad
 P_t=\theta_t^*L^{\boxplus t}.
\]
Then $P_t$ is ample and
\begin{equation}\label{eq:diag-epsilon}
 \varepsilon(P_t)\le\frac{(n+1)t^2}{b_0}.
\end{equation}
\end{proposition}

The proof compares two orders of vanishing of one determinant.
If \eqref{eq:diag-epsilon} fails, ordinary jets, specialized to large
positive characteristic, produce a column map of full generic rank.
Its rank drops on the diagonal, forcing a maximal minor to vanish
there to high order. A movable curve test, chosen in characteristic
zero from the order bound on $V$, gives a smaller upper bound for
that same minor. The next three estimates make both
bounds quantitative.

\subsection{Truncated powers, flags, and orders}

For a vector bundle $E$ on a smooth projective curve, let
$\mu_{\min}(E)$ and $\mu_{\max}(E)$ denote the extreme slopes in
its Harder--Narasimhan filtration. In characteristic $p>0$, let
$A_\bullet(E)$ be the symmetric algebra modulo the ideal generated
locally by $p$-th powers of degree-one elements. Thus in a frame of
a rank-$\rho$ bundle,
\[
 A_\bullet(E)=\OO[v_1,\ldots,v_\rho]/(v_1^p,\ldots,v_\rho^p).
\]
This definition is independent of the frame.

The relevant local model is the fibre product of $t$ copies of a
smooth $n$-fold $W$ over its relative Frobenius $F:W\to W'$, where
$W'$ is the Frobenius twist. In \'etale coordinates, differences
from the first factor give $n(t-1)$ variables whose $p$th powers vanish.
The associated graded along the reduced diagonal is therefore the
truncated algebra of $(\Omega_W^1)^{\oplus(t-1)}$. Below we use
the slope estimate for these graded pieces, the flag estimate to
bound their global sections, and the external-product estimate to
bound a determinant's order.

\begin{lemma}[Truncated-power slopes]\label{lem:diag-slopes}
Let $C$ be a smooth projective curve of genus $g$ over an
algebraically closed field of characteristic $p>0$. If $E$ has rank
$\rho$ and $\mu_{\min}(E)\ge0$, then, with
\[
 G=\max\{0,2g-2\},\qquad
 u=\rho(p-1),\qquad c_2=\rho(\rho-1)G,
\]
every nonzero graded piece satisfies
\begin{equation}\label{eq:diag-truncated-slope}
 \mu_{\max}(A_j(E))\le(p-1)\deg E+c_2,
 \qquad 0\le j\le u.
\end{equation}
The sum of their ranks is $p^\rho$.
\end{lemma}

\begin{proof}
We give the canonical-connection argument underlying the Frobenius
instability estimates of Langer \cite[Section~2]{Langer2004}; in
particular, we do not assume that tensor products are semistable in
characteristic $p$. For any rank-$\rho$ bundle $U$, absolute
Frobenius $F_C$ satisfies
\begin{equation}\label{eq:diag-one-step}
 \mu_{\min}(F_C^*U)
 \ge p\mu_{\min}(U)-(\rho-1)G.
\end{equation}
Indeed, consider the successive slopes of $F_C^*U$. If a successive
gap is larger than $G$, take the last such gap and let $S$ be the
preceding saturated filtration step. The second fundamental map of
the canonical connection,
\[
 S\longrightarrow ((F_C^*U)/S)\otimes\Omega_C^1,
\]
vanishes by the slope inequality. Consequently $S$ descends under
Frobenius. Here this descent can be checked directly: over $k(C)$,
put a matrix for $S$ in a chart with an identity minor. Differentiating
its columns gives vectors in $S\otimes\Omega_{k(C)/k}^1$ with zero
identity rows, hence all matrix entries have derivative zero. They
belong to $k(C)^p$. Extracting their roots and saturating gives a
subbundle $S_0\subset U$ with $S=F_C^*S_0$, since Frobenius is flat.

The quotient $Q=(F_C^*U)/S$ is therefore the pullback of a quotient
of $U$, so $\mu(Q)\ge p\mu_{\min}(U)$. All successive slope gaps
of $Q$ are at most $G$; its smallest slope is at least its average
minus $(\rho-1)G$. It is also the smallest slope of $F_C^*U$.
This proves \eqref{eq:diag-one-step}. If there is no large gap, use
$S=0$ and the same average-slope argument.

Iteration gives
\begin{equation}\label{eq:diag-iterated-slope}
 \mu_{\min}((F_C^e)^*E)
 \ge-(\rho-1)G\frac{p^e-1}{p-1}.
\end{equation}
Choose a line bundle $T_e$ of integer degree $d_e$ in the interval
\[
 2g-\mu_{\min}((F_C^e)^*E)
 \le d_e<2g+1-\mu_{\min}((F_C^e)^*E).
\]
The bundle $E_e=(F_C^e)^*E\otimes T_e$ is globally generated:
after subtracting one point, its minimum slope is at least $2g-1$;
Serre duality gives vanishing of $H^1$ and hence surjectivity of
evaluation. If $Q_i$ is a positive-rank torsion-free quotient of
$E^{\otimes i}$, then $(F_C^e)^*Q_i\otimes T_e^{\otimes i}$
is globally generated. Thus $p^e\mu(Q_i)+i d_e\ge0$.
Divide by $p^e$ and let $e$ tend to infinity, using
\eqref{eq:diag-iterated-slope}, to obtain
\begin{equation}\label{eq:diag-tensor-lower}
 \mu(Q_i)\ge-\frac{i(\rho-1)G}{p-1}.
\end{equation}
Since $A_i(E)$ is a quotient of $E^{\otimes i}$, its minimum slope
is at least $-c_2$ for $0\le i\le u$.

Multiplication gives a perfect pairing
$A_j(E)\otimes A_{u-j}(E)\to A_u(E)$: a truncated monomial pairs
with the unique complementary monomial. The top line is
$A_u(E)=(\det E)^{p-1}$, as is checked on diagonal and elementary
changes of frame. Therefore
\[
 A_j(E)^*\otimes(\det E)^{p-1}\simeq A_{u-j}(E).
\]
Its minimum slope is at least $-c_2$, proving
\eqref{eq:diag-truncated-slope}. Counting the $p^\rho$ monomials
with exponents between $0$ and $p-1$ gives the rank assertion.
\end{proof}

\begin{lemma}[Sections from a fixed flag]\label{lem:diag-flag}
Let $W$ be a smooth projective $n$-fold, let $H$ be an ample
rational Cartier divisor, and fix a smooth complete intersection
flag cut successively by divisors in $|l_iH|$, $1\le i<n$, ending
in a smooth integral curve $C$. Suppose that the $l_iH$ are integral
ample divisors. Set $d_i=l_iH\cdot C$. If a vector bundle $E$ of
rank $e$ satisfies $\mu_{\max}(E|_C)\le M$ with $M\ge0$, then
\begin{equation}\label{eq:diag-flag-bound}
 h^0(W,E)\le e(M+1)\prod_{i=1}^{n-1}(1+M/d_i).
\end{equation}
\end{lemma}

\begin{proof}
The divisor exact sequences along the flag bound $h^0(W,E)$ by
\[
 \sum_{k_1,\ldots,k_{n-1}\ge0}
 h^0\left(C,E|_C\otimes
       \OO_C\left(-\sum_{i=1}^{n-1}k_i l_iH\right)\right).
\]
Each iteration terminates because a sufficiently negative ample
twist of a fixed vector bundle has no sections. One may see this
by embedding that bundle in a finite sum of sufficiently positive
line bundles. A summand vanishes if $\sum_i k_i d_i>M$, by its
negative maximum slope. Otherwise evaluation at $\lfloor M\rfloor+1$
distinct points is injective, again by the slope criterion, so the
summand is at most $e(M+1)$. There are at most
$\prod_i(1+M/d_i)$ possible tuples. The same proof for $n=1$
uses an empty product.
\end{proof}

\begin{lemma}[Orders in an external product]\label{lem:diag-orders}
Let $W_i$ be integral projective varieties, $x_i$ smooth points,
and $Q_i$ line bundles with nonzero section spaces, for $1\le i\le t$.
Set $a_i=\max_{0\ne s\in H^0(W_i,Q_i)}\ord_{x_i}s$.
Every nonzero section of $\bigotimes_i\operatorname{pr}_i^*Q_i$
has order at most $\sum_i a_i$ at $(x_1,\ldots,x_t)$.
\end{lemma}

\begin{proof}
Choose bases adapted to the order filtrations of the section spaces.
Initial forms with the same order are linearly independent. By
K\"unneth, a section of the external product is a linear combination
of tensor products of these bases. In its smallest occurring total
degree, terms with different slot multidegrees cannot cancel; within
a multidegree their initial forms are independent because the slot
parameters are disjoint. Its order is therefore the smallest total
order of a basis tensor with nonzero coefficient, at most
$\sum_i a_i$.
\end{proof}

\subsection{The Frobenius rank comparison}

\begin{proof}[Proof of Proposition~\ref{prop:diagonal}]
The map $\theta_t$ is finite, so $P_t$ is ample. Fix $V,L,r,t$;
all auxiliary data in this proof may depend on these fixed choices.
Put
\begin{equation}\label{eq:diag-bq}
 b=b_0/t,\qquad q=2tr.
\end{equation}
Suppose that $\varepsilon(P_t)>(n+1)t/b$.

\medskip\noindent\emph{A fixed supply of ordinary jets.}
Choose a smooth very general $z\in Z_t$ whose coordinates lie over
$V_{\mathrm{reg}}$, and a rational number $c$ with
$(n+1)t<c<b\varepsilon(P_t;z)$. On the blow-up of $z$, the
pullback of $bP_t$ minus $c$ times the exceptional divisor is ample.
Indeed interpolate a nef class with slightly larger exceptional
coefficient and an ample class with small positive coefficient.
Serre vanishing on this blow-up, together with the usual pushforward
of powers of its exceptional ideal, gives, for sufficiently divisible
$k_0$, surjectivity onto ordinary jets through degree $(n+1)t k_0$.
Increase $k_0$ so that $Q=k_0bL$ is Cartier. Thus
\begin{equation}\label{eq:diag-fixed-jets}
 H^0(Z_t,k_0bP_t)\longrightarrow
 \OO_{Z_t,z}/\mathfrak m_z^{(n+1)t k_0+1}
 \quad\text{is surjective},
\end{equation}
where a local frame is understood.

Fix a trivialization $\eta$ of $rK_V$. On $V_{\mathrm{reg}}$,
the index cover is the torsor of canonical frames $\tau$ satisfying
$\tau^r=\eta$. The map $Y^t\to Z_t$ is etale over this locus.
At a chosen lift of $z$, K\"unneth and the character decomposition
express the pulled-back sections in \eqref{eq:diag-fixed-jets}
as combinations of tensors with slot factors
\begin{equation}\label{eq:diag-torsor-sections}
 \tau^{-m_i}s_i,\qquad
 s_i\in H^0(V_{\mathrm{reg}},Q+m_iK_V),\qquad
 0\le m_i<r,\qquad \sum_i m_i\equiv0\pmod r.
\end{equation}
Choose a smooth projective resolution $u:W\to V$ which is an
isomorphism over $V_{\mathrm{reg}}$. Write
$K_W=u^*K_V+A$, with $A\ge0$ exceptional. We use $L,Q$ also
for their pullbacks. Each $s_i$ extends to $Q+m_iK_W$: close its
effective divisor on $V$, pull back this rational Cartier divisor,
and add $m_iA$. Likewise
\[
 \eta_W\in H^0(W,rK_W),\qquad \Div(\eta_W)=rA.
\]
Fix finitely many tensors \eqref{eq:diag-torsor-sections} spanning
the required jets. Only this finite jet data will be specialized.

\medskip\noindent\emph{A fixed flag and a movable curve test.}
Since $K_WL^{n-1}=0$, we may take an ample rational divisor
$H=L+\delta H_{\mathrm{amp}}$, for $H_{\mathrm{amp}}$ fixed ample
and $\delta>0$ small rational, such that
\begin{equation}\label{eq:diag-H}
 \begin{split}
 H^n&\le3,\qquad LH^{n-1}\le H^n,\\
 (q+t-1)K_WH^{n-1}&\le b_0H^n.
 \end{split}
\end{equation}
The canonical divisor of $W$ is pseudo-effective. Generic
semipositivity \cite[Theorem~2.1]{CampanaPaun2015}, followed by
restriction of its Harder--Narasimhan factors
\cite[Theorem~6.1]{MehtaRamanathan1982}, gives a smooth complete
intersection flag in multiples $l_iH$ ending in a curve $C$ with
\begin{equation}\label{eq:diag-cotangent}
 \mu_{\min}(\Omega_W^1|_C)\ge0.
\end{equation}
The flag may be chosen to avoid the codimension-two singular sets
of the filtration quotients.

Independently choose a very general $x\in W$, in the isomorphism
locus of $u$, where the all-degree order bound defining $\gamma(L)$
holds. Let $\sigma:\widehat W\to W$ blow up $x$, with exceptional
divisor $J$, and define
\[
 D^+=bL+(q+1)K_W,\qquad d_0=b(C_1+1).
\]
In divisible degrees the sections of $D^+$ are those of $bL$
multiplied by a fixed exceptional section: $rK_V\sim0$, and
$u_*\OO_W(E)=\OO_V$ for effective integral exceptional $E$.
That fixed section has order zero at $x$. Hence the normalized
orders of sections of $D^+$ at $x$ are at most $bC_1$.

It follows that $\sigma^*D^+-d_0J$ is not pseudo-effective.
Otherwise interpolation with the big divisor $\sigma^*D^+$ would
make $\sigma^*D^+-dJ$ big for a rational $d$ strictly between
$bC_1$ and $d_0$, producing an excluded section. The projective
movable-cone duality theorem \cite[Theorem~2.2, preprint]{BDPP2013}
therefore gives a strongly movable curve class $\beta$ on
$\widehat W$ satisfying
\begin{equation}\label{eq:diag-test}
 D^+\cdot\beta<d_0J\cdot\beta,
 \qquad J\cdot\beta>0,
 \qquad L\cdot\beta\ge0,
 \qquad K_W\cdot\beta\ge0.
\end{equation}
Divisors from $W$ are pulled back in these pairings. More precisely,
we may choose $\beta=v_*(A_1\cdots A_{n-1})$, where
$v:T\to\widehat W$ is a smooth projective birational model and
the $A_i$ are very ample. The strict separation inequality holds on
one such generator of the movable cone. Nonnegative pairing with
$L$ and $K_W$ then forces $J\cdot\beta>0$.

\medskip\noindent\emph{Specialization of the finite data.}
Spread all the preceding algebraic data over an integral finitely
generated $\Z$-subalgebra of $\C$. This includes the points, blow-up,
flag, testing morphism, ample divisors, torsor, its marked lift,
and the finitely many jet-spanning sections. After shrinking the
base, preserve smoothness, geometric integrality, ampleness,
dominance of $v$, and the intersection numbers in
\eqref{eq:diag-H} and \eqref{eq:diag-test}. The fixed jet evaluation
is a map of finite locally free modules, so its surjectivity persists.
Invert $r$ and the finitely many denominators in the chosen data.

We can also preserve \eqref{eq:diag-cotangent}. A fixed relative
twist makes $\Omega_W^1|_C$ globally generated; thus the degrees of
its positive-rank quotients are bounded below uniformly by rank.
There are only finitely many negative degrees and ranks to exclude.
For each associated Hilbert polynomial, the relative Quot scheme
is projective and has empty geometric generic fibre, by
\eqref{eq:diag-cotangent}. Removing their images excludes all
negative-degree quotients. Quotients with torsion cause no issue,
since their torsion-free quotients have no larger degree.
Consequently \eqref{eq:diag-cotangent} holds on every remaining
geometric fibre.

There are such fibres in arbitrarily large characteristics $p$.
Work over their algebraically closed residue fields, retaining the
notation. The test $\beta$ still pairs nonnegatively with every
effective divisor: pull back by the retained dominant morphism $v$
and intersect with the very ample $A_i$. We have specialized a
fixed test, not an assertion about the entire pseudo-effective cone.

\medskip\noindent\emph{Full rank at a general product point.}
Let $F:W\to W'$ be relative Frobenius, where the prime indicates
the base-field twist. Set
\begin{equation}\label{eq:diag-ranks}
 N=\lfloor p/k_0\rfloor,\quad B=NQ,\quad
 \mathcal E=F_*\OO_W(B),\quad s=p^n,\quad R=s^t.
\end{equation}
The finite flat map $F$ has degree $s$, so $\mathcal E$ has rank $s$.
Products of $N$ of the fixed jet sections span jets through degree
$(n+1)t k_0N$. To check this, factor each monomial of degree at
most that number into $N$ monomials of degree at most
$(n+1)t k_0$, and choose corresponding sections with these initial
forms. Their products give a triangular spanning set for the jet
filtration. For large $p$,
\[
 (n+1)t k_0N\ge nt(p-1).
\]
Thus these products span the thick point defined by the $p$-th powers
of all $nt$ smooth parameters. This is the elementary comparison
of ordinary and Frobenius jets used in
\cite[Proposition~2.12]{MustataSchwede2014}. Etaleness identifies
the completed local rings at the torsor lift and its tuple in $W^t$,
and therefore identifies these thick points.

For an $N$-fold product of pure tensors, write $m_i'$ for the new
exponent in slot $i$. Then $0\le m_i'\le N(r-1)<pr$ and
$\sum_i m_i'\equiv0\pmod r$. Choose $a_i$ in $[r,2r-1]$ with
$pa_i\equiv m_i'\pmod r$. Their sum is a multiple of $r$ less
than $q=2tr$. Increase $a_1$ by the difference to obtain
\[
 a_i\ge0,\qquad \sum_i a_i=q,\qquad pa_i\ge m_i'.
\]
Multiplying the slot section by
$\eta_W^{(pa_i-m_i')/r}$ expresses its torsor factor as
$\tau^{-pa_i}$ times a section of $B+pa_iK_W$.
The frame factor is a nonzero constant on the thick Frobenius
point, because it is a $p$-th power of a unit.

By projection formula, the resulting column map is
\begin{equation}\label{eq:diag-column-map}
 \begin{split}
 \bigoplus_{\substack{a_i\ge0\\\sum_i a_i=q}}
 \left(\bigotimes_{i=1}^t H^0(W,B+pa_iK_W)\right)
 \otimes\OO_{(W')^t}\left(-\sum_{i=1}^t a_iK_{W',i}\right)
 \longrightarrow\mathcal E^{\boxtimes t}.
 \end{split}
\end{equation}
Here $K_{W',i}$ is pulled back from slot $i$.
The spanning assertion proves that this map has rank $R$ at one
point, and hence at the generic point. The identity used is the
line-bundle identity $F^*\OO_{W'}(K_{W'})=\OO_W(pK_W)$;
it is not the pullback map on differential forms.

\medskip\noindent\emph{Rank loss along the diagonal.}
On the diagonal $W'\subset(W')^t$, every source line bundle in
\eqref{eq:diag-column-map} restricts to $\OO_{W'}(-qK_{W'})$.
Let
\[
 \Delta_F=\underbrace{W\times_{W'}\cdots\times_{W'}W}_{t\text{ factors}},
 \qquad h:\Delta_F\to W',\qquad g:\Delta_F\to W
\]
be the structure map and the first projection. Finite base change
and projection formula identify
\[
 h_*\OO_{\Delta_F}\left(\sum_i\operatorname{pr}_i^*B
              +pq\operatorname{pr}_1^*K_W\right)
 = (\mathcal E^{\boxtimes t})|_{W'}\otimes\OO_{W'}(qK_{W'}).
\]
After twisting, the columns restrict to global sections of this
bundle, so their rank at every diagonal point is bounded by the
dimension of that global section space.

Filter by the ideal of the reduced diagonal in $\Delta_F$ and
push forward by $g$. The graded pieces are
\begin{equation}\label{eq:diag-graded}
 \mathcal G_j=\OO_W(tB+pqK_W)
          \otimes A_j((\Omega_W^1)^{\oplus(t-1)}),
 \quad 0\le j\le n(t-1)(p-1).
\end{equation}
Indeed, in etale coordinates the $t-1$ difference slots have
$n(t-1)$ variables with $p$-th powers zero. Their degree-one
transitions are those of the cotangent bundle. This is the
many-factor form of the canonical Frobenius-diagonal filtration
\cite[Theorem~3.7]{Sun2008},
\cite[Definition~3.1 and Corollary~3.5]{KitadaiSumihiro2008}.
If $e_j$ is the rank of the truncated factor, then
$\sum_j e_j=p^{n(t-1)}$.

Apply Lemma~\ref{lem:diag-slopes} on $C$ with
$\rho=n(t-1)$ and its constant $c_2$, independent of $p,j$.
Put $D=l_1\cdots l_{n-1}H^n$, so $l_iH\cdot C=l_iD$.
Using $B\cdot C\le pbL\cdot C$, \eqref{eq:diag-H}, and
$K_W\cdot C\ge0$, we obtain
\begin{align*}
 \mu_{\max}(\mathcal G_j|_C)
 &\le \bigl(tB+pqK_W+(p-1)(t-1)K_W\bigr)\cdot C+c_2\\
 &\le M_p:=2b_0pD+c_2.
\end{align*}
Lemma~\ref{lem:diag-flag} now gives
\begin{equation}\label{eq:diag-piece-bound}
 \begin{split}
 h^0(W,\mathcal G_j)
 &\le e_j(M_p+1)\prod_{i=1}^{n-1}(1+M_p/(l_iD))\\
 &=e_jp^n\bigl((2b_0)^nH^n+O(p^{-1})\bigr).
 \end{split}
\end{equation}
The error is uniform in $j$: all the flag data and $c_2$ were
fixed before varying $p$. The leading coefficient follows from
$D^n/\prod_i(l_iD)=H^n$.
Summing the filtration bounds, and using \eqref{eq:diag-b0}, gives
\begin{equation}\label{eq:diag-rank-loss}
 \operatorname{rank}\bigl(\text{column map at a diagonal point}\bigr)
 \le p^{nt}\bigl((2b_0)^nH^n+O(p^{-1})\bigr)
 \le R/4
\end{equation}
for sufficiently large $p$.

We have obtained the two rank bounds. It remains to apply the
characteristic-zero curve test to a nonzero maximal minor, after
specialization, and compare its order with this diagonal rank loss.

\medskip\noindent\emph{The determinant contradiction.}
Choose $R$ columns with nonzero determinant. If $a_i^{(\nu)}$ is
the slot-$i$ weight in column $\nu$, their determinant is a section
of the external product of the line bundles
\[
 \mathcal Q_i=(\det\mathcal E)^{\otimes R/s}
       \otimes\OO_{W'}\left(\sum_{\nu=1}^R a_i^{(\nu)}K_{W'}\right).
\]
Set $\lambda_i=R^{-1}\sum_\nu a_i^{(\nu)}$, so $0\le\lambda_i\le q$.
Identifying numerical classes under the base-field twist,
\begin{equation}\label{eq:diag-determinant-class}
 \frac{c_1(\mathcal Q_i)}R
 =\frac Bp+\left(\frac{p-1}{2p}+\lambda_i\right)K_W.
\end{equation}
For a direct verification, the two-factor version of
\eqref{eq:diag-graded} filters $F^*\mathcal E$ by
$\OO_W(B)\otimes A_j(\Omega_W^1)$. Among all $p^n=s$
truncated monomials, the total exponent of each variable is
$s(p-1)/2$. Hence
\[
 c_1(F^*\mathcal E)=sB+\frac{s(p-1)}2K_W.
\]
Frobenius pullback multiplies numerical divisor classes by $p$,
which proves \eqref{eq:diag-determinant-class} and agrees with
\cite[Lemma~4.2]{Sun2008}.

The normalized class in \eqref{eq:diag-determinant-class} pairs
with the specialized test $\beta$ to at most $D^+\cdot\beta$,
because $B/p$ is a nonnegative scalar multiple of $L$ no larger
than $bL$, and the coefficient of $K_W$ is at most $q+1$.
For any nonzero slot section $s_i$ of $\mathcal Q_i$, set
$m_i=\ord_{x'}s_i$. Its divisor has effective strict transform
of class $\sigma^*c_1(\mathcal Q_i)-m_iJ'$ on the twisted
blow-up. Thus \eqref{eq:diag-test} gives
\[
 m_i(J\cdot\beta)
 \le c_1(\mathcal Q_i)\cdot\beta
 \le R D^+\cdot\beta<Rd_0(J\cdot\beta),
\]
and $m_i<Rd_0$.
The determinant is nonzero, so all its slot section spaces are
nonzero. Lemma~\ref{lem:diag-orders} gives the strict upper bound
for its order at $(x',\ldots,x')$:
\[
 Rt d_0=Rb_0(C_1+1)<3R/4.
\]
On the other hand, \eqref{eq:diag-rank-loss} lets us make at
least $R-\lfloor R/4\rfloor$ rows of its matrix vanish at that
point, by constant invertible row operations in local frames.
Every determinant term then lies in the corresponding power of
the maximal ideal, giving order at least $3R/4$. The contradiction
proves \eqref{eq:diag-epsilon}.
\end{proof}

\section{Compatible leaves and rational transport}\label{sec:transport}

We now use the quadratic growth in Proposition~\ref{prop:diagonal}
to finish the high-index contradiction. Low-cost curves on many
diagonal products generate compatible leaves. Large covering degree
forces these leaves to project generically finitely onto their images
in every coordinate, so their dimensions stabilize. Their projection degrees cannot grow
exponentially under a polynomial bound. A degree-one forgetting
map then supplies rational evaluation of one coordinate from the
others, and turns local flows into birational self-maps.
This is the many-factor argument of \cite[Section~8]{U4}, in
arbitrary fixed dimension.

\begin{proposition}[The high-index contradiction]\label{prop:index-contradiction}
There is no sequence of finite surjective covers
$\pi_\nu:Y_\nu\to V_\nu$ of normal projective complex $n$-folds,
generically torsors for cyclic groups $G_\nu$ of orders
$r_\nu\to\infty$, satisfying the
following conditions. The groups $\Bir(V_\nu)$ are countable,
and there are ample rational Cartier divisors $L_\nu$ and constants
$c,C_2>0$, independent of $\nu$, such that
\[
 \varepsilon(L_\nu)\ge c,\qquad
 \varepsilon(\pi_\nu^*L_\nu)\ge cr_\nu^{1/n}.
\]
For $t\ge2$, let $Z_{\nu,t}=Y_\nu^t/G_{\nu,\mathrm{diag}}$,
let $\theta_{\nu,t}:Z_{\nu,t}\to V_\nu^t$ be the finite
quotient map, and set $P_{\nu,t}=\theta_{\nu,t}^*L_\nu^{\boxplus t}$.
The final condition is
\[
 \varepsilon(P_{\nu,t})\le C_2t^2
\]
for every fixed $t$, on a tail allowed to depend on $t$.
All Seshadri constants are at very general points.
\end{proposition}

\begin{proof}
Fix an integer $T$, eventually large in terms of $n,c,C_2$.
On a tail of the sequence, use the upper bounds for all
$2\le t\le T$ simultaneously. Suppress $\nu$ until the covering
degree must vary, and put $B=C_2T^2+1$.
Since $G$ is abelian, $\theta_t$ is the finite quotient by
$G^t/G_{\mathrm{diag}}$: its further quotient has function field
$\C(V^t)$, and normality identifies it with $V^t$.

\medskip\noindent\emph{Curves and projection-compatible leaves.}
For each $t$, the strict bound $\varepsilon(P_t)<B$ supplies a
marked covering family of integral curves with
$P_t$-degree divided by marked multiplicity at most $B$.
Choose a curve at geometric generic data and spread it with its
mark. Equivalently, define the finite data over a countable field
and realize the resulting generic families over $\C$.
Include every deck translate, every permutation of the slots,
and every nonconstant image family under ordered coordinate
projections to smaller sizes between $2$ and $T$. Closing under
these operations still gives finite lists; no uniform bound on
their sizes is needed.

The cost cannot increase under projection. Indeed, for a
nonconstant map $D\to D'$ of integral curves of degree $e$,
with marked point $z$ over $z'$, a general local hyperplane
through $z'$ gives, on normalization,
\[
 e\operatorname{mult}_{z'}D'\ge\operatorname{mult}_zD.
\]
The contribution from branches through $z$ cannot exceed the
total over $z'$, and is at least $\operatorname{mult}_zD$.
Since $P_t$ minus the pullback of the lower polarization is nef,
the degree inequality and this multiplicity inequality give the
asserted cost comparison. The projected marks remain dominant;
we take reduced integral geometric generic image curves.

Apply Lemma~\ref{lem:chain-input} to these lists.
Let $H_t$ be a geometric generic leaf and $h_t=\dim H_t$.
Then
\begin{equation}\label{eq:transport-chain-degree-new}
 h_t>0,\qquad P_t^{h_t}\cdot H_t\le(h_tB)^{h_t}.
\end{equation}
Its finite-quotient functoriality shows that
$W_t=\theta_t(H_t)$ is the geometric generic leaf
of a rational quotient
$\xi_t:V^t\dashrightarrow S_t$ with geometrically integral
generic fibre. These quotients and distributions are invariant
under permutations of the slots.

The image of an upper leaf under a coordinate projection lies
in the corresponding lower leaf. To check the genericity in
this assertion, retain the parameter of every upper family when
forming its image family. Equality of the values of the lower quotient
at a mark and a generic endpoint is an identity on that generic
incidence. It remains valid with the additional history of any
upper chain. Every upper step therefore preserves that quotient value;
a stationary projected step does so immediately. This proves
the inclusion for generic chains and then for their leaf closures.
All comparisons below use a generic tuple and its projections,
with their quotient parameter fields extended to a common
algebraically closed field.

\medskip\noindent\emph{Finiteness of the single-slot projections.}
For fixed $T$ and sufficiently large $r$, each $H_t$ projects
generically finitely onto its image in every $V$-slot. Suppose
otherwise. For one slot write $I$ for the image and
$j=\dim I$, $a=h_t-j>0$. The marked point of $I$ is ambient
very general and generic on $I$ over the leaf's quotient parameter
field. Thus
\begin{equation}\label{eq:transport-image-bound-new}
 L^j\cdot I\ge c^j.
\end{equation}
Here and below the degree bound follows by blowing up the marked
point and taking the top intersection of the nef class given
by the Seshadri bound on the strict transform of the subvariety.
The generic point is smooth on the subvariety and has multiplicity
one. For $j=0$ the assertion reads $1\ge1$.

Over a geometric generic point of this slot, the fibre of
$Z_t\to V$ identifies with $Y^{t-1}$: choose one lift in the
torsor to fix the first coordinate. The polarization restricts
to the external sum of the remaining copies of $\pi^*L$.
Projecting a curve through the marked tuple to a nonconstant
slot, using the preceding degree--multiplicity comparison,
shows that this restricted polarization has Seshadri constant
at least $cr^{1/n}$ there. For an $a$-dimensional component $J$
of the geometric generic fibre of $H_t\to I$ through the mark,
we obtain $P_t^a\cdot J\ge(cr^{1/n})^a$.
The projection formula and nefness of
$P_t-\operatorname{pr}^*L$ now give
\[
 P_t^{h_t}\cdot H_t
 \ge(\operatorname{pr}^*L)^jP_t^a\cdot H_t
 \ge c^j(cr^{1/n})^a.
\]
This contradicts \eqref{eq:transport-chain-degree-new} as
$r\to\infty$, since $h_t\le nT$ and $T$ is fixed.
Only finitely many slots and sizes are involved, so the conclusion
holds simultaneously on a tail. In particular,
\begin{equation}\label{eq:transport-dimension-new}
 1\le h_t\le n,
\end{equation}
and the same single-slot finiteness holds for $W_t$.

\medskip\noindent\emph{A degree-one forgetting map.}
Let $d_t$ be the degree of $W_t$ over its first-slot image.
The degree of $H_t$ over that image contains $d_t$ as a factor.
Combining \eqref{eq:transport-chain-degree-new} with
\eqref{eq:transport-image-bound-new} gives
\begin{equation}\label{eq:transport-degree-new}
 d_t\le KT^{2n},\qquad
 K=\max_{1\le h\le n}\left(\frac{h(C_2+1)}c\right)^h.
\end{equation}
The constant $K$ is independent of $T$.
Projection compatibility and finite first-slot projection imply
$h_t\le h_{t-1}$ for $t\ge3$. Whenever equality holds, the
upper leaf maps onto the integral lower leaf; their first-slot
images coincide and
\begin{equation}\label{eq:transport-ratio-new}
 d_t/d_{t-1}=\deg(W_t\longrightarrow W_{t-1})\in\N.
\end{equation}
There is a constant stretch in the $T-1$ dimensions of length
at least $(T-1)/n$. Choose $T$ so large that
\[
 (T-1)/n-1>\log_2(KT^{2n}).
\]
If all the adjacent degrees on that stretch were at least two,
\eqref{eq:transport-ratio-new} would contradict
\eqref{eq:transport-degree-new}. Hence for some $k\ge3$,
\begin{equation}\label{eq:transport-degree-one-new}
 h_k=h_{k-1}=h>0,\qquad
 W_k\longrightarrow W_{k-1}\ \text{is birational}.
\end{equation}
Permutation invariance gives the same conclusion for every
single-slot omission.

We now have a positive-dimensional leaf which determines each
missing coordinate rationally. The rest of the proof turns this
rational recovery into too many birational transformations of $V$.

\medskip\noindent\emph{Transport of tangent directions.}
The tangent distribution $\mathcal D_k=\ker(d\xi_k)$ has rank
$h$ at a generic tuple. It projects injectively to each single-slot
tangent space, and isomorphically to the lower distribution under
every omission. These assertions follow from separability and the
generically finite maps of geometric generic leaves.
Let $A_i\subset T_{V,x_i}$ be its slot-$i$ image. It depends
rationally only on $x_i$. Indeed, in a rational Grassmann chart
its coordinates belong to the function field of every
$(k-1)$-tuple containing slot $i$, so they belong to the
intersection of these fields, namely $\C(V)$ in slot $i$.
Likewise, the transport from $A_i$ to $A_j$, obtained by lifting
to $\mathcal D_k$ and projecting to slot $j$, belongs to the
field of the pair $(x_i,x_j)$.

The field intersection just used is elementary for independent
product variables. Replace an omitted variable by a fresh
independent generic variable. A rational function pulled from
the remaining slots is unchanged by this replacement. Repeating
and then specializing the fresh variables to general constants
shows that it is pulled from the common slots. Applied to the
Grassmann and matrix chart coordinates, this proves both descent
statements, including when $k=3$.

Permutation symmetry therefore gives one rational plane
distribution $A(x)$ and rational isomorphisms
$R(x,y):A(x)\to A(y)$. Lifting a vector to the same leaf tangent
space and using three slots gives
\begin{equation}\label{eq:transport-cocycle-new}
 R(y,z)R(x,y)=R(x,z).
\end{equation}
Choose a general reference point $a\in V(\C)$ and transport a
basis of $A(a)$ by $R(a,x)$. This produces rational vector fields
$v_1,\ldots,v_h$, independent at a generic point, whose simultaneous
copies $(v_i,\ldots,v_i)$ span $\mathcal D_k$ by
\eqref{eq:transport-cocycle-new}. No commutativity is asserted
or needed.

\medskip\noindent\emph{Rational recovery and local flows.}
Write $\xi=\xi_k$ and let $D$ be the image closure of
\[
 \Lambda=(\operatorname{pr}_{<k},\xi):V^k\dashrightarrow D.
\]
This map is birational. To see that no algebraic degree remains,
write $F_0\subset E\subset F$ for the function fields of
$S_k,D,V^k$. The field $E$ is generated over $F_0$ by the first
$k-1$ slots, and \eqref{eq:transport-degree-one-new} shows that
$F/E$ is finite. Geometric integrality of the generic leaf means
that $F/F_0$ is regular. After extension to an algebraic closure
of $F_0$, the corresponding generic fields therefore stay fields,
and the degree of $F/E$ remains unchanged. The geometric
forgetting map has degree one, so $F=E$.
Let $\operatorname{ev}$ be the last coordinate of the rational
inverse of $\Lambda$.

Choose a general complex tuple $(u,y_0)$, with $u\in V^{k-1}$,
in smooth opens where the vector fields, quotient, and rational
evaluation are regular. For small $z=(z_1,\ldots,z_h)\in\C^h$,
let $E_z$ be the ordered composition of the time-$z_i$ local
holomorphic flows of $v_i$. Their coefficients are holomorphic
on the chosen opens, so the local existence and uniqueness
theorem for differential equations gives the flows uniformly
on smaller neighbourhoods. The reversed composition of the
negative-time flows is their local inverse $E_z^-$.

The simultaneous flows preserve $\xi$. Consequently, for $y$
near $y_0$ and sufficiently small $z$,
\begin{equation}\label{eq:transport-flow-new}
 E_z(y)=\operatorname{ev}\bigl(E_z(u),\xi(u,y)\bigr).
\end{equation}
For each fixed $z$ the right side is rational in $y$. Equality
on an analytic open proves that its input lies in $D$ and meets
the domain of evaluation; that analytic open is Zariski dense.
Thus it defines a rational self-map of $V$, regular at $y_0$
after shrinking the times. The identical argument for $E_z^-$
gives a rational inverse, because the two compositions agree
with the identity on smaller analytic opens. Each $E_z$ is
therefore birational. Since $h>0$ and the $v_i(y_0)$ are
independent, the map $z\mapsto E_z(y_0)$ has nonzero derivative
and uncountable image. This produces uncountably many distinct
members of $\Bir(V)$, a contradiction.
\end{proof}

Apply this proposition to the sequence in
Proposition~\ref{prop:index-reduction}. Choose $L$ rationally
scaled so that $1\le L^n\le2$. Proposition~\ref{prop:scalar}
gives uniform $\gamma(L)\le C_1$, $\varepsilon(L)\ge c$, and
$\varepsilon(\pi^*L)\ge cr^{1/n}$, since
$(\pi^*L)^n=rL^n$. Proposition~\ref{prop:diagonal} supplies the
remaining bound with $C_2=(n+1)/b_0$, while countability of
$\Bir(V)$ is part of Proposition~\ref{prop:index-reduction}.
The high-index sequence is impossible. Together with the other
cases of that reduction, this proves the zero-boundary klt
canonical-index bound in dimension $n$.
